%% kapa: exam=1 round=32 year=2021 period=1 form=A subject=economics subject_ko=경제학원론 seq=38 number=78
%% source: Q-Net 감정평가사 기출문제 (cst00309) artlSeq=5208397 file=extracted/1008/5208397/A형/2021년도 제32회 감정평가사 1차 1교시 A형.pdf
%% conversion: kordoc 4.24.51 -> tools/gen1.py (auto); review=proofread-sonnet; flags=none
\begin{question}{78}
A국 경제의 총수요곡선과 총공급곡선이 각각 $P=-Y_d+4$, $P=P_e+(Y_s-2)$이다. $P_e$가 3에서 5로 증가할 때, (ㄱ)\uline{균형소득수준}과 (ㄴ)\uline{균형물가수준}의 변화는? (단, $P$는 물가수준, $Y_d$는 총수요, $Y_s$는 총공급, $P_e$는 기대물가수준이다.)
\begin{choices}
  \item ㄱ: 상승, ㄴ: 상승
  \item ㄱ: 하락, ㄴ: 상승
  \item ㄱ: 상승, ㄴ: 하락
  \item ㄱ: 하락, ㄴ: 하락
  \item ㄱ: 불변, ㄴ: 불변
\end{choices}
\end{question}
\answer{2}
